\hypertarget{objetivos}{%
\section{Objetivos}\label{objetivos}}

\begin{frame}{Objetivos}
\protect\hypertarget{objetivos-1}{}
\begin{itemize}
\item
  Apresentar a linguagem de especifica\c{c}\~{a}o do Happy.
\item
  Apresentar exemplos de analisadores LALR construídos utilizando o
  Happy.
\end{itemize}
\end{frame}

\hypertarget{introduuxe7uxe3o-ao-happy}{%
\section{Introdu\c{c}\~{a}o ao Happy}\label{introduuxe7uxe3o-ao-happy}}

\begin{frame}{Motiva\c{c}\~{a}o}
\protect\hypertarget{motivauxe7uxe3o}{}
\begin{itemize}
\tightlist
\item
  Construir tabelas LALR \`{a} m\~{a}o é impraticável para gramáticas reais
\item
  Assim como o \textbf{Alex} automatiza análise léxica, o \textbf{Happy}
  automatiza análise sintática
\end{itemize}
\end{frame}

\begin{frame}[fragile]{Motiva\c{c}\~{a}o}
\protect\hypertarget{motivauxe7uxe3o-1}{}
\begin{itemize}
\tightlist
\item
  Happy gera analisadores \textbf{LALR(1)} a partir de especifica\c{c}\~{o}es
  gramaticais no estilo BNF
\item
  O próprio compilador \textbf{GHC} usa um parser gerado pelo Happy
\item
  Análogo ao \passthrough{\lstinline!yacc!} e
  \passthrough{\lstinline!Bison!} do ecossistema C/C++
\end{itemize}
\end{frame}

\begin{frame}[fragile]{Estrutura de um Arquivo
\passthrough{\lstinline!.y!}}
\protect\hypertarget{estrutura-de-um-arquivo-.y}{}
\begin{lstlisting}
{ cabe\c{c}alho Haskell }

diretivas e declara\c{c}\~{a}o de tokens

%%

regras gramaticais

{ rodapé Haskell }
\end{lstlisting}
\end{frame}

\begin{frame}[fragile]{Estrutura de um Arquivo
\passthrough{\lstinline!.y!}}
\protect\hypertarget{estrutura-de-um-arquivo-.y-1}{}
\begin{itemize}
\tightlist
\item
  Quatro partes separadas pelo marcador \passthrough{\lstinline!\%\%!}:
\end{itemize}

\begin{enumerate}
\tightlist
\item
  \textbf{Cabe\c{c}alho}: módulo e imports Haskell (copiado literalmente)
\item
  \textbf{Diretivas}: configura\c{c}\~{o}es do parser e declara\c{c}\~{a}o de tokens
\item
  \textbf{Regras}: produ\c{c}\~{o}es BNF com a\c{c}\~{o}es sem\^{a}nticas em Haskell
\item
  \textbf{Rodapé}: fun\c{c}\~{o}es auxiliares
\end{enumerate}
\end{frame}

\begin{frame}[fragile]{Gerando o Parser}
\protect\hypertarget{gerando-o-parser}{}
\begin{lstlisting}[language=bash]
happy ExpParser.y            # gera ExpParser.hs
happy -i ExpParser.y         # também gera ExpParser.info (tabela)
happy -o ExpParser.hs ExpParser.y
\end{lstlisting}
\end{frame}

\begin{frame}[fragile]{Gerando o Parser}
\protect\hypertarget{gerando-o-parser-1}{}
\begin{itemize}
\tightlist
\item
  Com Cabal: basta listar \passthrough{\lstinline!happy!} em
  \passthrough{\lstinline!build-tool-depends!}
\end{itemize}
\end{frame}

\hypertarget{diretivas-do-happy}{%
\section{Diretivas do Happy}\label{diretivas-do-happy}}

\begin{frame}[fragile]{Principais Diretivas}
\protect\hypertarget{principais-diretivas}{}
\begin{longtable}[]{@{}ll@{}}
\toprule\noalign{}
Diretiva & Significado \\
\midrule\noalign{}
\endhead
\passthrough{\lstinline!\%name p N!} & Gera fun\c{c}\~{a}o
\passthrough{\lstinline!p!} iniciando pelo n\~{a}o-terminal
\passthrough{\lstinline!N!} \\
\bottomrule\noalign{}
\end{longtable}
\end{frame}

\begin{frame}[fragile]{Principais Diretivas}
\protect\hypertarget{principais-diretivas-1}{}
\begin{longtable}[]{@{}ll@{}}
\toprule\noalign{}
Diretiva & Significado \\
\midrule\noalign{}
\endhead
\passthrough{\lstinline!\%tokentype \{ T \}!} & Define o tipo Haskell
dos tokens consumidos \\
\bottomrule\noalign{}
\end{longtable}
\end{frame}

\begin{frame}[fragile]{Principais Diretivas}
\protect\hypertarget{principais-diretivas-2}{}
\begin{longtable}[]{@{}
  >{\raggedright\arraybackslash}p{(\columnwidth - 2\tabcolsep) * \real{0.4714}}
  >{\raggedright\arraybackslash}p{(\columnwidth - 2\tabcolsep) * \real{0.5286}}@{}}
\toprule\noalign{}
\begin{minipage}[b]{\linewidth}\raggedright
Diretiva
\end{minipage} & \begin{minipage}[b]{\linewidth}\raggedright
Significado
\end{minipage} \\
\midrule\noalign{}
\endhead
\passthrough{\lstinline!\%monad \{ M \}\{ (>>=) \}\{ return \}!} &
Executa o parser dentro da m\^{o}nada \passthrough{\lstinline!M!} \\
\bottomrule\noalign{}
\end{longtable}
\end{frame}

\begin{frame}[fragile]{Principais Diretivas}
\protect\hypertarget{principais-diretivas-3}{}
\begin{longtable}[]{@{}ll@{}}
\toprule\noalign{}
Diretiva & Significado \\
\midrule\noalign{}
\endhead
\passthrough{\lstinline!\%lexer \{ f \}\{ eof \}!} & Lexer incremental:
\passthrough{\lstinline!f!} obtém o próximo token \\
\bottomrule\noalign{}
\end{longtable}
\end{frame}

\begin{frame}[fragile]{Principais Diretivas}
\protect\hypertarget{principais-diretivas-4}{}
\begin{longtable}[]{@{}ll@{}}
\toprule\noalign{}
Diretiva & Significado \\
\midrule\noalign{}
\endhead
\passthrough{\lstinline!\%error \{ h \}!} & Chama
\passthrough{\lstinline!h!} ao encontrar erro de sintaxe \\
\bottomrule\noalign{}
\end{longtable}
\end{frame}

\begin{frame}[fragile]{Principais Diretivas}
\protect\hypertarget{principais-diretivas-5}{}
\begin{longtable}[]{@{}
  >{\raggedright\arraybackslash}p{(\columnwidth - 2\tabcolsep) * \real{0.1667}}
  >{\raggedright\arraybackslash}p{(\columnwidth - 2\tabcolsep) * \real{0.8333}}@{}}
\toprule\noalign{}
\begin{minipage}[b]{\linewidth}\raggedright
Diretiva
\end{minipage} & \begin{minipage}[b]{\linewidth}\raggedright
Significado
\end{minipage} \\
\midrule\noalign{}
\endhead
\passthrough{\lstinline!\%expect n!} & Declara exatamente
\passthrough{\lstinline!n!} conflitos shift-reduce esperados \\
\bottomrule\noalign{}
\end{longtable}
\end{frame}

\begin{frame}[fragile]{Principais Diretivas}
\protect\hypertarget{principais-diretivas-6}{}
\begin{longtable}[]{@{}ll@{}}
\toprule\noalign{}
Diretiva & Significado \\
\midrule\noalign{}
\endhead
\passthrough{\lstinline!\%left!}, \passthrough{\lstinline!\%right!},
\passthrough{\lstinline!\%nonassoc!} & Preced\^{e}ncia e associatividade \\
\bottomrule\noalign{}
\end{longtable}
\end{frame}

\begin{frame}[fragile]{Diretivas}
\protect\hypertarget{diretivas}{}
\begin{lstlisting}
%name parser Exp
%monad {Alex}{(>>=)}{return}
%tokentype { Token }
%error     { parseError }
%lexer {lexer}{Token _ TEOF}
\end{lstlisting}
\end{frame}

\begin{frame}[fragile]{Diretivas}
\protect\hypertarget{diretivas-1}{}
\begin{itemize}
\tightlist
\item
  \passthrough{\lstinline!\%monad!} integra com a m\^{o}nada
  \passthrough{\lstinline!Alex!} do lexer gerado pelo Alex
\item
  \passthrough{\lstinline!\%lexer!} faz Happy chamar
  \passthrough{\lstinline!lexer!} para obter cada token (modo
  incremental)
\item
  O modo incremental é necessário quando o lexer é monádico
\end{itemize}
\end{frame}

\begin{frame}[fragile]{Declara\c{c}\~{a}o de Tokens}
\protect\hypertarget{declarauxe7uxe3o-de-tokens}{}
\begin{itemize}
\tightlist
\item
  A se\c{c}\~{a}o \passthrough{\lstinline!\%token!} associa nomes simbólicos a
  padr\~{o}es de correspond\^{e}ncia:
\end{itemize}

\begin{lstlisting}
%token
      num       {Token _ (TNumber $$)}
      '('       {Token _ TLParen}
      ')'       {Token _ TRParen}
      '+'       {Token _ TPlus}
      '*'       {Token _ TTimes}
\end{lstlisting}
\end{frame}

\begin{frame}[fragile]{Declara\c{c}\~{a}o de Tokens}
\protect\hypertarget{declarauxe7uxe3o-de-tokens-1}{}
\begin{itemize}
\tightlist
\item
  \passthrough{\lstinline!$$!} em
  \passthrough{\lstinline!Token \_ (TNumber $$)!} captura o valor
  inteiro
\end{itemize}

--- Disponível nas a\c{c}\~{o}es como \passthrough{\lstinline!$1!},
\passthrough{\lstinline!$2!}, etc.

\begin{itemize}
\tightlist
\item
  Tokens sem valor (como \passthrough{\lstinline!TLParen!}) apenas
  verificam o construtor
\end{itemize}
\end{frame}

\begin{frame}[fragile]{Declara\c{c}\~{o}es de Preced\^{e}ncia}
\protect\hypertarget{declarauxe7uxf5es-de-preceduxeancia}{}
\begin{lstlisting}
%left '+'
%left '*'
\end{lstlisting}
\end{frame}

\begin{frame}[fragile]{Declara\c{c}\~{o}es de Preced\^{e}ncia}
\protect\hypertarget{declarauxe7uxf5es-de-preceduxeancia-1}{}
\begin{itemize}
\tightlist
\item
  A \textbf{ordem} das linhas determina preced\^{e}ncia: de menor para maior
\item
  \passthrough{\lstinline!*!} aparece depois de
  \passthrough{\lstinline!+!} → \passthrough{\lstinline!*!} tem
  \textbf{maior} preced\^{e}ncia
\end{itemize}
\end{frame}

\begin{frame}[fragile]{Declara\c{c}\~{o}es de Preced\^{e}ncia}
\protect\hypertarget{declarauxe7uxf5es-de-preceduxeancia-2}{}
\begin{itemize}
\tightlist
\item
  \passthrough{\lstinline!\%left!}: associatividade \`{a} esquerda
\item
  \passthrough{\lstinline!\%right!}: associatividade \`{a} direita
\item
  \passthrough{\lstinline!\%nonassoc!}: n\~{a}o associativo (erro se
  encadeado)
\end{itemize}
\end{frame}

\hypertarget{regras-gramaticais-e-auxe7uxf5es}{%
\section{Regras Gramaticais e
A\c{c}\~{o}es}\label{regras-gramaticais-e-auxe7uxf5es}}

\begin{frame}[fragile]{Regras Gramaticais}
\protect\hypertarget{regras-gramaticais}{}
\begin{itemize}
\tightlist
\item
  Após o \passthrough{\lstinline!\%\%!}, produ\c{c}\~{o}es BNF com a\c{c}\~{o}es em
  Haskell:
\end{itemize}

\begin{lstlisting}
Exp : num         { EInt $1 }
    | Exp '+' Exp { $1 :+: $3 }
    | Exp '*' Exp { $1 :*: $3 }
    | '(' Exp ')' { $2 }
\end{lstlisting}
\end{frame}

\begin{frame}[fragile]{Regras Gramaticais}
\protect\hypertarget{regras-gramaticais-1}{}
\begin{itemize}
\tightlist
\item
  \passthrough{\lstinline!$1!}, \passthrough{\lstinline!$2!},
  \passthrough{\lstinline!$3!}: valores sem\^{a}nticos dos símbolos (da
  esquerda para a direita)
\item
  A\c{c}\~{a}o executada quando o parser \textbf{reduz} por aquela produ\c{c}\~{a}o
\item
  Tipo inferido pelo compilador a partir das a\c{c}\~{o}es:
  \passthrough{\lstinline!Exp!}
\end{itemize}
\end{frame}

\begin{frame}[fragile]{Rodapé: Fun\c{c}\~{o}es Auxiliares}
\protect\hypertarget{rodapuxe9-funuxe7uxf5es-auxiliares}{}
\begin{lstlisting}[language=Haskell]
{
lexer :: (Token -> Alex a) -> Alex a
lexer = (=<< alexMonadScan)

expParser :: String -> IO (Either String Exp)
expParser content = pure $ runAlex content parser
}
\end{lstlisting}
\end{frame}

\hypertarget{gerando-a-tabela-lalr}{%
\section{Gerando a Tabela LALR}\label{gerando-a-tabela-lalr}}

\begin{frame}[fragile]{A Op\c{c}\~{a}o \passthrough{\lstinline!-i!}}
\protect\hypertarget{a-opuxe7uxe3o--i}{}
\begin{lstlisting}[language=bash]
happy -i ExpParser.y
# gera ExpParser.info
\end{lstlisting}
\end{frame}

\begin{frame}[fragile]{Arquivos .info}
\protect\hypertarget{arquivos-.info}{}
O arquivo \passthrough{\lstinline!.info!} contém:

\begin{enumerate}
\tightlist
\item
  \textbf{Gramática aumentada} com produ\c{c}\~{o}es numeradas
\item
  \textbf{Conjuntos FIRST} de cada n\~{a}o-terminal
\end{enumerate}
\end{frame}

\begin{frame}{Arquivos .info}
\protect\hypertarget{arquivos-.info-1}{}
\begin{enumerate}
\setcounter{enumi}{2}
\tightlist
\item
  \textbf{Estados do aut\^{o}mato LALR}:cada um com seus itens LR(1) e
  lookaheads
\item
  \textbf{Tabela de a\c{c}\~{o}es}: shift, reduce, accept para cada (estado,
  terminal)
\end{enumerate}
\end{frame}

\begin{frame}{Arquivos .info}
\protect\hypertarget{arquivos-.info-2}{}
\begin{enumerate}
\setcounter{enumi}{4}
\tightlist
\item
  \textbf{Tabela goto}: para cada (estado, n\~{a}o-terminal)
\item
  \textbf{Conflitos detectados}: com estado e token envolvidos
\end{enumerate}
\end{frame}

\begin{frame}[fragile]{Exemplo de Estado no
\passthrough{\lstinline!.info!}}
\protect\hypertarget{exemplo-de-estado-no-.info}{}
\begin{lstlisting}
State 7

        Exp -> Exp . '+' Exp          (rule 2)
        Exp -> Exp . '*' Exp          (rule 3)
        Exp -> Exp '+' Exp .          (rule 2)

        '*'    shift, and enter state 5
               (reduce using rule 2)

        '+'    shift, and enter state 4
               (reduce using rule 2)

        ')'    reduce using rule 2
        '$'    reduce using rule 2
\end{lstlisting}
\end{frame}

\hypertarget{conflitos}{%
\section{Conflitos}\label{conflitos}}

\begin{frame}[fragile]{Shift-Reduce}
\protect\hypertarget{shift-reduce}{}
\begin{itemize}
\tightlist
\item
  Ocorre quando o parser pode tanto fazer shift quanto reduce no mesmo
  estado e token.
\end{itemize}

\textbf{Causa comum}: gramática de express\~{o}es escrita de forma ambígua:

\begin{lstlisting}
Exp : num
    | Exp '+' Exp
    | Exp '*' Exp
    | '(' Exp ')'
\end{lstlisting}
\end{frame}

\begin{frame}[fragile]{Shift-Reduce}
\protect\hypertarget{shift-reduce-1}{}
\begin{itemize}
\tightlist
\item
  Para \passthrough{\lstinline!1 + 2 * 3!}, após empilhar
  \passthrough{\lstinline!Exp '+' Exp!} (com
  \passthrough{\lstinline!Exp!} = \passthrough{\lstinline!2!}) e ver
  \passthrough{\lstinline!*!}:

  \begin{itemize}
  \tightlist
  \item
    \textbf{Reduce} \passthrough{\lstinline!Exp '+' Exp!} → interpreta
    \passthrough{\lstinline!(1+2)*3!}
  \item
    \textbf{Shift} \passthrough{\lstinline!*!} → interpreta
    \passthrough{\lstinline!1+(2*3)!} (correto!)
  \end{itemize}
\item
  Happy retorna:
  \passthrough{\lstinline!ExpParser.y: shift/reduce conflicts: 4!}
\end{itemize}
\end{frame}

\begin{frame}{Solu\c{c}\~{a}o: Preced\^{e}ncia}
\protect\hypertarget{soluuxe7uxe3o-preceduxeancia}{}
Happy usa as declara\c{c}\~{o}es para resolver:

\begin{itemize}
\tightlist
\item
  Token com \textbf{maior} preced\^{e}ncia que a produ\c{c}\~{a}o → \textbf{shift}
\item
  Token com \textbf{menor} preced\^{e}ncia → \textbf{reduce}
\item
  Token com \textbf{mesma} preced\^{e}ncia → aplica associatividade
\end{itemize}
\end{frame}

\begin{frame}[fragile]{Shift-Reduce}
\protect\hypertarget{shift-reduce-2}{}
\begin{itemize}
\tightlist
\item
  Gramática para if/else
\end{itemize}

\begin{lstlisting}
Stmt : 'if' Exp 'then' Stmt
     | 'if' Exp 'then' Stmt 'else' Stmt
     | 'other'
\end{lstlisting}
\end{frame}

\begin{frame}[fragile]{Shift-Reduce}
\protect\hypertarget{shift-reduce-3}{}
\begin{itemize}
\tightlist
\item
  Para \passthrough{\lstinline!if e1 then if e2 then s1 else s2!}:

  \begin{itemize}
  \tightlist
  \item
    Shift \passthrough{\lstinline!else!} → pertence ao
    \passthrough{\lstinline!if!} mais próximo (interno): conven\c{c}\~{a}o
    padr\~{a}o
  \item
    Reduce \passthrough{\lstinline!if ... then s1!} →
    \passthrough{\lstinline!else!} pertence ao
    \passthrough{\lstinline!if!} externo
  \end{itemize}
\item
  O comportamento padr\~{a}o do Happy (preferir shift) resolve esse problema
  corretamente.
\end{itemize}
\end{frame}

\begin{frame}[fragile]{Reduce-Reduce}
\protect\hypertarget{reduce-reduce}{}
\begin{itemize}
\tightlist
\item
  Ocorre quando dois itens completos diferentes reivindicam a mesma
  entrada.
\end{itemize}

\begin{lstlisting}
Prog : Decl
     | Expr

Decl : 'id'
Expr : 'id'
\end{lstlisting}
\end{frame}

\begin{frame}[fragile]{Reduce-Reduce}
\protect\hypertarget{reduce-reduce-1}{}
Ao ver \passthrough{\lstinline!id!} e precisar reduzir:

\begin{itemize}
\tightlist
\item
  N\~{a}o sabe se deve criar \passthrough{\lstinline!Decl!} ou
  \passthrough{\lstinline!Expr!}
\item
  Happy escolhe a \textbf{primeira regra} (aviso, n\~{a}o erro)
\end{itemize}
\end{frame}

\begin{frame}{Reduce-Reduce}
\protect\hypertarget{reduce-reduce-2}{}
\begin{itemize}
\tightlist
\item
  Diferentemente do shift-reduce, conflitos reduce-reduce raramente t\^{e}m
  resolu\c{c}\~{a}o padr\~{a}o correta.
\end{itemize}

--- \textbf{Quase sempre indicam uma gramática problemática}.
\end{frame}

\begin{frame}[fragile]{Reduce-Reduce}
\protect\hypertarget{reduce-reduce-3}{}
\begin{itemize}
\tightlist
\item
  Solu\c{c}\~{a}o: Tornar os dois n\~{a}o-terminais sintáticos distintos por
  contexto:
\end{itemize}

\begin{lstlisting}
-- Antes: ambíguo
Prog : 'id'   -- pode ser Decl ou Expr?

-- Depois: discriminar por palavra-chave
Prog : 'let' 'id'    -- declara\c{c}\~{a}o
     | 'id'          -- express\~{a}o
\end{lstlisting}
\end{frame}

\begin{frame}[fragile]{Reduce-Reduce}
\protect\hypertarget{reduce-reduce-4}{}
\begin{itemize}
\tightlist
\item
  Outra abordagem: remover ambiguidade na análise sem\^{a}ntica
\end{itemize}

--- Exemplo: \passthrough{\lstinline!typedef!} de C, onde
identificadores de tipo requerem consulta ao contexto.
\end{frame}

\hypertarget{conclusuxe3o}{%
\section{Conclus\~{a}o}\label{conclusuxe3o}}

\begin{frame}{Conclus\~{a}o}
\protect\hypertarget{conclusuxe3o-1}{}
\begin{itemize}
\item
  Apresentamos a linguagem do gerador de analisadores sintáticos LALR
  Happy.
\item
  Apresentamos um exemplo de uma gramática de express\~{o}es.
\end{itemize}
\end{frame}
